\chapter{Cédric Villani --- Fields Medal 2010}\label{ch2010cedricvillani:cuxe9dric-villani-fields-medal-2010-a-comprehensive-report-on-his-work-and-impact}

\begin{center}\rule{0.5\linewidth}{0.5pt}\end{center}

\textbf{Prepared:} July 2026 \textbf{Subject:} Cédric Villani (b.
October 5, 1973, Brive-la-Gaillarde, France) \textbf{Position:}
Professor, Institut Henri Poincaré (Director, 2009--2017); Professor,
Université Lyon 1; Member of the French Parliament (2017--2022)
\textbf{Doctoral Advisor:} Pierre-Louis Lions \textbf{Fields Medal
Citation:} \emph{``For his proofs of nonlinear Landau damping and
convergence to equilibrium for the Boltzmann equation.''}

\begin{center}\rule{0.5\linewidth}{0.5pt}\end{center}

\section{Table of Contents}\label{ch2010cedricvillani:table-of-contents}

\begin{enumerate}
\def\labelenumi{\arabic{enumi}.}
\tightlist
\item
 Introduction
\item
 Biographical Background and Education
\item
 Kinetic Theory: Background and Context
\item
 The Boltzmann Equation
\item
 The H-Theorem and Convergence to Equilibrium
\item
 Villani's Work on the Boltzmann Equation
\item
 The Cercignani Conjecture
\item
 Optimal Transport: Background
\item
 The Wasserstein Metric and the Monge--Ampère Equation
\item
 Villani's Books on Optimal Transport
\item
 Nonlinear Landau Damping
\item
 The Mouhot--Villani Proof of Landau Damping
\item
 The Landau Equation and Plasma Physics
\item
 Hypocoercivity
\item
 Entropy and Information Theory
\item
 Geometric Methods in Functional Inequalities
\item
 Public Engagement and Science Policy
\item
 Awards and Recognition
\item
 Selected Major Publications
\item
 Ongoing Work and Future Directions
\item
 Conclusion
\item
 References and Further Reading
\end{enumerate}

\begin{prereqs} Measure theory (probability measures, weak convergence), basic PDE theory (transport equations), functional analysis. \end{prereqs}

\begin{objectives} Explain the optimal transport framework and its applications; describe the Landau damping result; understand the entropy method for PDEs. \end{objectives}


\begin{center}\rule{0.5\linewidth}{0.5pt}\end{center}

\section{1. Introduction}\label{ch2010cedricvillani:introduction}

Cédric Villani was awarded the Fields Medal in 2010 for his fundamental
contributions to kinetic theory and optimal transport. His work is
characterized by a rare combination of deep analytical power, geometric
intuition, and a willingness to tackle the most difficult problems in
mathematical physics.

Villani's two Fields Medal--worthy achievements are: 1.
\textbf{Convergence to equilibrium for the Boltzmann equation:} He
proved that solutions of the Boltzmann equation converge to equilibrium
at an exponential rate, resolving a long-standing conjecture. 2.
\textbf{Nonlinear Landau damping:} With Clément Mouhot, he proved that
the Landau damping phenomenon --- the exponential decay of plasma waves
in a collisionless plasma --- is a rigorous mathematical consequence of
the Vlasov equation.

Beyond these specific results, Villani has made substantial
contributions to optimal transport theory --- a field that was already
important but that Villani helped to develop into a central tool in
analysis, geometry, and probability. His two books on optimal transport
have become the standard references in the field.

Villani is also notable for his public engagement with mathematics. As
Director of the Institut Henri Poincaré (2009--2017) and later as a
member of the French Parliament (2017--2022), he has been one of the
most visible mathematicians in the public sphere.

\begin{center}\rule{0.5\linewidth}{0.5pt}\end{center}

\section{2. Biographical Background and
Education}\label{ch2010cedricvillani:biographical-background-and-education}

\subsection{2.1 Early Life}\label{ch2010cedricvillani:early-life}

Cédric Villani was born on October 5, 1973, in Brive-la-Gaillarde, in
the Corrèze department of central France. He was raised in a family of
teachers: his father was a mathematics teacher and his mother was an
English teacher and later a school principal.

Villani showed mathematical talent early but was also interested in
literature and music. He has described his intellectual formation as
shaped by both the rigor of mathematics and the creativity of the arts.

\subsection{2.2 Education}\label{ch2010cedricvillani:education}

Villani attended the prestigious Lycée Louis-le-Grand in Paris, one of
France's most selective secondary schools. He then entered the École
Normale Supérieure (ENS) in Paris, the most elite institution of higher
education in France, in 1992.

\subsection{2.3 PhD Studies}\label{ch2010cedricvillani:phd-studies}

Villani completed his PhD at Université Paris-Dauphine in 1998 under the
supervision of Pierre-Louis Lions --- himself a Fields Medalist (1994).
His thesis, ``Contribution à l'étude mathématique des équations de
Boltzmann et de Landau en théorie cinétique des gaz,'' already showed
his deep interest in kinetic theory.

\subsection{2.4 Academic Career}\label{ch2010cedricvillani:academic-career}

After his PhD, Villani held positions at: - \textbf{CNRS:} Researcher
(1998--2000) - \textbf{Université Paris-Dauphine:} Professor
(2000--2008) - \textbf{Institut Henri Poincaré:} Director (2009--2017) -
\textbf{Université Lyon 1:} Professor (2008--present)

\subsection{2.5 Political Career}\label{ch2010cedricvillani:political-career}

In 2017, Villani was elected to the French National Assembly as a member
of La République En Marche!, representing the Essonne department. He
served until 2022, focusing on science policy, innovation, and
education. His time in politics attracted international attention, with
profiles in publications ranging from \emph{The New York Times} to
\emph{Le Monde}.

\begin{center}\rule{0.5\linewidth}{0.5pt}\end{center}

\section{3. Kinetic Theory: Background and
Context}\label{ch2010cedricvillani:kinetic-theory-background-and-context}

\subsection{3.1 What Is Kinetic Theory?}\label{ch2010cedricvillani:what-is-kinetic-theory}

Kinetic theory describes the behavior of a gas at the
\textbf{mesoscopic} scale --- between the microscopic scale (individual
atoms) and the macroscopic scale (fluid dynamics). Instead of tracking
individual particles (impossible for Avogadro's number of molecules),
kinetic theory uses a \textbf{distribution function} \(f(t, x, v)\) that
gives the density of particles at position \(x\) with velocity \(v\) at
time \(t\).

\subsection{3.2 The Key Equations}\label{ch2010cedricvillani:the-key-equations}

The fundamental equations of kinetic theory are: - \textbf{Boltzmann
equation:} Describes rarefied gases with binary collisions -
\textbf{Vlasov equation:} Describes collisionless plasmas -
\textbf{Landau equation:} Describes plasmas with Coulomb collisions -
\textbf{Fokker--Planck equation:} Describes particles in a random medium

\subsection{3.3 The Challenge}\label{ch2010cedricvillani:the-challenge}

Kinetic equations pose fundamental mathematical challenges: -
\textbf{Existence and uniqueness:} Do solutions exist for all time? -
\textbf{Regularity:} Are solutions smooth? - \textbf{Long-time
behavior:} Do solutions approach equilibrium? How fast? -
\textbf{Hydrodynamic limit:} Do kinetic equations reduce to fluid
equations in the appropriate limit?

Villani made major contributions to all of these questions.

\begin{center}\rule{0.5\linewidth}{0.5pt}\end{center}

\section{4. The Boltzmann Equation}\label{ch2010cedricvillani:the-boltzmann-equation}

\subsection{4.1 The Equation}\label{ch2010cedricvillani:the-equation}

The Boltzmann equation, introduced by Ludwig Boltzmann in 1872,
describes the evolution of a dilute gas:

\[
\frac{\partial f}{\partial t} + v \cdot \nabla_x f = Q(f, f)
\]

where \(Q(f, f)\) is the \textbf{collision operator}:

\[
Q(f, f)(v) = \int_{\mathbb{R}^3} \int_{S^2} B(|v - v_*|, \theta) [f(v') f(v'_*) - f(v) f(v_*)] \, d\sigma \, dv_*
\]

Here, \(v, v_*\) are pre-collision velocities, \(v', v'_*\) are
post-collision velocities, and \(B\) is the collision kernel depending
on the intermolecular potential.

\subsection{4.2 The H-Theorem}\label{ch2010cedricvillani:the-h-theorem}

Boltzmann's famous H-theorem states that the entropy:

\[
H(f) = \int f \log f \, dv \, dx
\]

is non-increasing along solutions of the Boltzmann equation:

\[
\frac{dH}{dt} \le 0
\]

Equality holds only when \(f\) is a Maxwellian (Gaussian in velocity):

\[
f_{\text{eq}}(v) = \frac{\rho}{(2\pi T)^{3/2}} e^{-|v - u|^2 / (2T)}
\]

\subsection{4.3 The Fundamental
Question}\label{ch2010cedricvillani:the-fundamental-question}

Boltzmann's H-theorem says entropy decreases, but it does not say
\textbf{how fast}. The question of the \textbf{rate of convergence to
equilibrium} --- whether entropy decays exponentially, polynomially, or
at some other rate --- was one of the central open problems in kinetic
theory for over a century.

\begin{center}\rule{0.5\linewidth}{0.5pt}\end{center}

\section{5. The H-Theorem and Convergence to
Equilibrium}\label{ch2010cedricvillani:the-h-theorem-and-convergence-to-equilibrium}

\subsection{5.1 The Entropy--Entropy Production
Method}\label{ch2010cedricvillani:the-entropyentropy-production-method}

A general approach to proving convergence to equilibrium is the
\textbf{entropy--entropy production method}. If one can prove an
inequality of the form:

\[
D(f) \ge \kappa \cdot (H(f) - H(f_{\text{eq}}))
\]

where \(D(f) = -dH/dt\) is the entropy production, then:

\[
\frac{d}{dt} (H - H_{\text{eq}}) \le -\kappa (H - H_{\text{eq}})
\]

which gives exponential decay:

\[
H(t) - H_{\text{eq}} \le (H(0) - H_{\text{eq}}) e^{-\kappa t}
\]

\subsection{5.2 The Difficulty}\label{ch2010cedricvillani:the-difficulty}

For the Boltzmann equation, proving such an inequality requires
understanding the spectral gap of the linearized collision operator, as
well as controlling the \textbf{nonlinear} contributions. The linearized
operator has a spectral gap (this was known from the work of Grad,
Caflisch, and others), but extending this to the full nonlinear equation
was extremely challenging.

\begin{center}\rule{0.5\linewidth}{0.5pt}\end{center}

\section{6. Villani's Work on the Boltzmann
Equation}\label{ch2010cedricvillani:villanis-work-on-the-boltzmann-equation}

\subsection{6.1 The Main Result}\label{ch2010cedricvillani:the-main-result}

In a series of papers (2002--2009), Villani, often in collaboration with
Laurent Desvillettes, proved:

\textbf{Theorem (Desvillettes--Villani):} For solutions of the Boltzmann
equation with various collision kernels (including hard spheres), the
convergence to equilibrium is exponentially fast, provided the initial
data is sufficiently regular and the domain is bounded with suitable
boundary conditions (e.g., a periodic box).

The rate of exponential convergence depends only on: - The collision
kernel - The geometry of the domain - Macroscopic parameters (density,
temperature)

\subsection{6.2 The Method}\label{ch2010cedricvillani:the-method}

The proof involves: 1. \textbf{Spectral analysis:} Understanding the
spectrum of the linearized Boltzmann operator 2.
\textbf{Entropy--entropy production estimates:} Proving a lower bound
for entropy production in terms of relative entropy 3.
\textbf{Regularity estimates:} Controlling the nonlinear terms using
smoothing properties of the collision operator

\subsection{6.3 Impact}\label{ch2010cedricvillani:impact}

This result resolved a fundamental question that had been open since
Boltzmann's H-theorem (1872). It showed that the H-theorem is not just a
statement about monotonic decay, but about exponential decay --- the gas
approaches equilibrium at a definite, predictable rate.

\begin{center}\rule{0.5\linewidth}{0.5pt}\end{center}

\section{7. The Cercignani Conjecture}\label{ch2010cedricvillani:the-cercignani-conjecture}

\subsection{7.1 The Conjecture}\label{ch2010cedricvillani:the-conjecture}

In the 1980s, Carlo Cercignani conjectured a precise inequality relating
entropy and entropy production for the spatially homogeneous Boltzmann
equation:

\[
D(f) \ge \frac{\lambda}{2} H(f|f_{\text{eq}})
\]

where \(H(f|f_{\text{eq}}) = \int f \log(f/f_{\text{eq}})\) is the
relative entropy and \(\lambda\) is the spectral gap of the linearized
operator.

\subsection{7.2 Villani's Work on the
Conjecture}\label{ch2010cedricvillani:villanis-work-on-the-conjecture}

Villani made important contributions to the Cercignani conjecture. While
the full conjecture (with the optimal constant) remains open in general,
Villani and collaborators proved: - The inequality holds with a positive
constant (possibly smaller than \(\lambda/2\)) for many collision
kernels. - For Maxwellian molecules (a special case), the conjecture is
true with the optimal constant. - The conjecture is false in certain
regimes (the constant may depend on the solution).

\subsection{7.3 Significance}\label{ch2010cedricvillani:significance}

The Cercignani conjecture, if fully resolved, would provide the optimal
rate of convergence to equilibrium for the spatially homogeneous
Boltzmann equation. Villani's work clarified the status of the
conjecture and provided the best available results.

\begin{center}\rule{0.5\linewidth}{0.5pt}\end{center}

\section{8. Optimal Transport:
Background}\label{ch2010cedricvillani:optimal-transport-background}

\subsection{8.1 The Monge Problem}\label{ch2010cedricvillani:the-monge-problem}

Optimal transport began with Gaspard Monge's 1781 problem: given two
distributions of mass (e.g., a pile of sand and a target shape), find
the most efficient way to transport the mass from the source to the
target.

Monge's problem: Find a map \(T\) such that:

\[
\int |T(x) - x| \, d\mu(x) \text{ is minimized subject to } T_*\mu = \nu
\]

\subsection{8.2 The Kantorovich
Relaxation}\label{ch2010cedricvillani:the-kantorovich-relaxation}

In 1942, Leonid Kantorovich proposed a relaxed formulation: instead of a
map, consider a \textbf{transport plan} --- a joint distribution \(\pi\)
on the product space:

\[
\int |x - y| \, d\pi(x, y) \text{ minimized subject to } \pi \text{ having marginals } \mu, \nu
\]

\subsection{8.3 The Wasserstein
Distance}\label{ch2010cedricvillani:the-wasserstein-distance}

The \textbf{Wasserstein distance} measures the minimal cost of
transporting one distribution to another:

\[
W_p(\mu, \nu) = \left( \inf_{\pi} \int |x - y|^p \, d\pi(x, y) \right)^{1/p}
\]

This metric has become a fundamental tool in many areas of mathematics.

\begin{center}\rule{0.5\linewidth}{0.5pt}\end{center}

\section{9. The Wasserstein Metric and the Monge--Ampère
Equation}\label{ch2010cedricvillani:the-wasserstein-metric-and-the-mongeampuxe8re-equation}

\subsection{9.1 Optimal Transport and
PDE}\label{ch2010cedricvillani:optimal-transport-and-pde}

A significant connection, discovered by Brenier (1991) and further
developed by McCann, Otto, and others, links optimal transport to
partial differential equations. The optimal transport map satisfies the
\textbf{Monge--Ampère equation}:

\[
\det(\nabla^2 \varphi) \circ T^{-1} = \frac{\mu}{\nu \circ T^{-1}}
\]

\subsection{9.2 The Otto Calculus}\label{ch2010cedricvillani:the-otto-calculus}

Felix Otto (2001) realized that the Wasserstein metric endows the space
of probability measures with a formal Riemannian structure, and many
PDEs (the heat equation, the Fokker--Planck equation, etc.) can be
interpreted as \textbf{gradient flows} in this metric space.

\subsection{9.3 Villani's Contributions}\label{ch2010cedricvillani:villanis-contributions}

Villani made fundamental contributions to the theory of optimal
transport: - \textbf{Regularity of optimal transport maps:} With Figalli
and De Philippis, he proved that optimal transport maps are smooth under
natural conditions. - \textbf{Wasserstein geometry:} He developed the
geometric theory of Wasserstein spaces. - \textbf{Gradient flows:} He
studied the gradient flow structure of PDEs in Wasserstein space.

\begin{center}\rule{0.5\linewidth}{0.5pt}\end{center}

\section{10. Villani's Books on Optimal
Transport}\label{ch2010cedricvillani:villanis-books-on-optimal-transport}

\subsection{\texorpdfstring{10.1 \emph{Topics in Optimal Transportation}
(2003)}{10.1 Topics in Optimal Transportation (2003)}}\label{ch2010cedricvillani:topics-in-optimal-transportation-2003}

This book, based on a graduate course, provides a comprehensive
introduction to optimal transport theory. It covers: - The Monge and
Kantorovich problems - The Wasserstein metric - The Brenier theorem -
Connections to PDE

\subsection{\texorpdfstring{10.2 \emph{Optimal Transport: Old and New}
(2009)}{10.2 Optimal Transport: Old and New (2009)}}\label{ch2010cedricvillani:optimal-transport-old-and-new-2009}

This massive treatise (1000+ pages) is the definitive reference on
optimal transport. It covers: - The complete theory from Monge to the
present - Applications to geometry, analysis, and probability - Open
problems and future directions

\subsection{10.3 Impact}\label{ch2010cedricvillani:impact-1}

These books have been instrumental in popularizing optimal transport
theory and making it accessible to a wide range of mathematicians. They
are cited thousands of times and are considered essential references for
anyone working in the field.

\begin{center}\rule{0.5\linewidth}{0.5pt}\end{center}

\section{11. Nonlinear Landau Damping}\label{ch2010cedricvillani:nonlinear-landau-damping}

\subsection{11.1 The Physical Phenomenon}\label{ch2010cedricvillani:the-physical-phenomenon}

Landau damping, predicted by Lev Landau in 1946, is a phenomenon in
plasma physics where plasma waves (Langmuir waves) decay exponentially
in time, even in the absence of collisions. The damping is caused by the
interaction of the wave with particles moving at the wave's phase
velocity --- particles moving slightly faster than the wave give energy
to the wave, while particles moving slightly slower absorb energy. If
the velocity distribution has a negative slope at the wave speed, the
net effect is damping.

\subsection{11.2 The Linear Theory}\label{ch2010cedricvillani:the-linear-theory}

Landau's calculation was for the \textbf{linearized} Vlasov--Poisson
equation. He showed that the electric field decays as:

\[
E(t) \sim e^{-\gamma t}
\]

where \(\gamma\) is the \textbf{Landau damping rate}, determined by the
slope of the velocity distribution at the wave speed.

\subsection{11.3 The Nonlinear Question}\label{ch2010cedricvillani:the-nonlinear-question}

The crucial question: does Landau damping persist in the \textbf{full
nonlinear} Vlasov--Poisson equation? For decades, this was believed to
be true but was not rigorously proved. The difficulty is that nonlinear
effects could, in principle, produce lasting echoes or other phenomena
that prevent decay.

\begin{center}\rule{0.5\linewidth}{0.5pt}\end{center}

\section{12. The Mouhot--Villani Proof of Landau
Damping}\label{ch2010cedricvillani:the-mouhotvillani-proof-of-landau-damping}

\subsection{12.1 The Main Theorem}\label{ch2010cedricvillani:the-main-theorem}

In a 2010 paper, Clément Mouhot and Cédric Villani proved:

\textbf{Theorem:} For the Vlasov--Poisson equation on a periodic domain
\(\mathbb{T}^d \times \mathbb{R}^d\), if the initial distribution
function is sufficiently close to a stable homogeneous state (e.g., a
Maxwellian), then: 1. The distribution function converges to a**
homogeneous** state (a function of velocity only) as \(t \to \infty\).
2. The electric field decays exponentially to zero. 3. The decay rate is
given by the linear Landau damping rate (to leading order).

\subsection{12.2 The Method}\label{ch2010cedricvillani:the-method-1}

The proof uses: - \textbf{Hamiltonian mechanics:} The Vlasov--Poisson
system is a Hamiltonian system on the space of distribution functions. -
\textbf{Complex analysis:} The Laplace transform method, extending
Landau's original calculation. - \textbf{Nonlinear analysis:} A
Newton-type iteration scheme to control nonlinear interactions. -
\textbf{Return to equilibrium:} Use of the linear damping result in a
perturbative framework.

\subsection{12.3 Key Ingredients}\label{ch2010cedricvillani:key-ingredients}

The technical tools include: - \textbf{Analytic norms:} The use of
Gevrey regularity (intermediate between \(C^\infty\) and analytic) to
control the nonlinearity - \textbf{Dispersion estimates:} Quantitative
bounds on the transport of particles - \textbf{KAM-type ideas:} The
proof uses a ``normal form'' transformation to eliminate resonant
interactions

\subsection{12.4 Significance}\label{ch2010cedricvillani:significance-1}

The Mouhot--Villani proof resolved one of the important open
problems in kinetic theory. Landau damping had been a cornerstone of
plasma physics for over 60 years, but its rigorous mathematical
justification was missing. The proof showed that Landau damping is not
just a linear phenomenon but persists in the full nonlinear setting,
under appropriate conditions.

\begin{center}\rule{0.5\linewidth}{0.5pt}\end{center}

\section{13. The Landau Equation and Plasma
Physics}\label{ch2010cedricvillani:the-landau-equation-and-plasma-physics}

\subsection{13.1 The Landau Equation}\label{ch2010cedricvillani:the-landau-equation}

The Landau equation is the kinetic equation for a plasma with Coulomb
collisions (long-range interactions). It is similar in structure to the
Boltzmann equation but with a different collision operator (the Landau
collision operator).

\subsection{13.2 Villani's
Contributions}\label{ch2010cedricvillani:villanis-contributions-1}

Villani made fundamental contributions to the Landau equation: -
\textbf{Existence of weak solutions:} Proved global existence for the
spatially homogeneous Landau equation. - \textbf{Regularity:} Showed
that solutions become smooth instantaneously. - \textbf{Asymptotic
behavior:} Proved convergence to equilibrium.

\subsection{13.3 The Landau Equation vs.~the Boltzmann
Equation}\label{ch2010cedricvillani:the-landau-equation-vs.-the-boltzmann-equation}

The Landau equation is in some ways more singular than the Boltzmann
equation (the collision kernel has a strong angular singularity), but it
also has more smoothing properties. Villani exploited these properties
to obtain results that were not yet available for the Boltzmann
equation.

\begin{center}\rule{0.5\linewidth}{0.5pt}\end{center}

\section{14. Hypocoercivity}\label{ch2010cedricvillani:hypocoercivity}

\subsection{14.1 The Concept}\label{ch2010cedricvillani:the-concept}

Hypocoercivity is a concept introduced by Villani (in a 2009 monograph)
to describe the \textbf{return to equilibrium} for kinetic equations.
While a system may not be directly coercive (the dissipation may not act
on all degrees of freedom), the combination of transport and dissipation
can nevertheless produce exponential convergence.

\subsection{14.2 The Hypocoercivity
Framework}\label{ch2010cedricvillani:the-hypocoercivity-framework}

The key idea: Even if the collision operator only acts on the velocity
variable, the \textbf{transport term} \(v \cdot \nabla_x\) couples
position and velocity, eventually allowing the collisions to drive the
entire system to equilibrium.

Villani's hypocoercivity framework provides a systematic way to prove
exponential convergence for such systems.

\subsection{14.3 Applications}\label{ch2010cedricvillani:applications}

The hypocoercivity framework has been applied to: - The Boltzmann
equation - The Landau equation - The Fokker--Planck equation - Kinetic
models of sociodynamics - Kinetic models of chemotaxis

\begin{center}\rule{0.5\linewidth}{0.5pt}\end{center}

\section{15. Entropy and Information
Theory}\label{ch2010cedricvillani:entropy-and-information-theory}

\subsection{15.1 Entropy Methods}\label{ch2010cedricvillani:entropy-methods}

Throughout his work, Villani has used entropy as a central tool: -
\textbf{Entropy production:} Quantifying how fast entropy decreases -
\textbf{Relative entropy:} Measuring distance to equilibrium -
\textbf{Fisher information:} A refinement of entropy that gives sharper
estimates

\subsection{15.2 The H-Theorem Revisited}\label{ch2010cedricvillani:the-h-theorem-revisited}

Villani's work on the H-theorem went beyond Boltzmann's original
formulation: - \textbf{Quantitative H-theorem:} Not just that entropy
decreases, but how fast - \textbf{Optimal constants:} Determining the
best possible rates - \textbf{Nonlinear effects:} How nonlinear
interactions affect the entropy decay

\subsection{15.3 Connections to Information
Theory}\label{ch2010cedricvillani:connections-to-information-theory}

Villani has explored the connections between kinetic theory and
information theory: - The Boltzmann entropy is the same as the Shannon
entropy (up to a sign) - The Fisher information plays a central role in
both fields (through de Bruijn's identity in information theory and
through the entropy production in kinetic theory) - Optimal transport
provides a geometric interpretation of information-theoretic quantities

\begin{center}\rule{0.5\linewidth}{0.5pt}\end{center}

\section{16. Geometric Methods in Functional
Inequalities}\label{ch2010cedricvillani:geometric-methods-in-functional-inequalities}

\subsection{16.1 Functional Inequalities}\label{ch2010cedricvillani:functional-inequalities}

Functional inequalities --- such as Sobolev inequalities, Poincaré
inequalities, and logarithmic Sobolev inequalities --- are fundamental
tools in analysis. Villani used optimal transport to prove and sharpen
these inequalities.

\subsection{16.2 The Brunn--Minkowski
Inequality}\label{ch2010cedricvillani:the-brunnminkowski-inequality}

The Brunn--Minkowski inequality in convex geometry states that the
volume of the Minkowski sum of two sets satisfies certain concavity
properties. Villani (building on work of McCann) showed that this
inequality can be understood through optimal transport, and he proved
sharper versions.

\subsection{16.3 The Prékopa--Leindler
Inequality}\label{ch2010cedricvillani:the-pruxe9kopaleindler-inequality}

The Prékopa--Leindler inequality is a functional generalization of the
Brunn--Minkowski inequality. Villani used optimal transport to give a
new proof and to obtain the optimal constants.

\subsection{16.4 The Talagrand
Inequality}\label{ch2010cedricvillani:the-talagrand-inequality}

The Talagrand inequality relates the Wasserstein distance to the
relative entropy:

\[
W_2(\mu, \nu) \le \sqrt{2 H(\mu|\nu)}
\]

Villani extended this inequality to various settings (Riemannian
manifolds, metric measure spaces) and used it as a tool in his study of
the Boltzmann equation.

\begin{center}\rule{0.5\linewidth}{0.5pt}\end{center}

\section{17. Public Engagement and Science
Policy}\label{ch2010cedricvillani:public-engagement-and-science-policy}

\subsection{17.1 Director of the Institut Henri
Poincaré}\label{ch2010cedricvillani:director-of-the-institut-henri-poincaruxe9}

As director of IHP (2009--2017), Villani transformed the institute into
a major center for mathematical research and public outreach. He
established: - Monthly thematic programs - Public lecture series - The
``Mathematics for Everyone'' initiative - The IHP's role as a hub for
international collaborations

\subsection{17.2 Member of the French
Parliament}\label{ch2010cedricvillani:member-of-the-french-parliament}

In 2017, Villani was elected to the French National Assembly for the
Essonne department. During his term (2017--2022): - He authored a major
report on artificial intelligence (``AI for Humanity'') - He advocated
for increased research funding - He served on the Parliamentary Office
for Scientific and Technological Assessment - He wrote a book on the
role of mathematics in society

\subsection{17.3 Public Intellectual}\label{ch2010cedricvillani:public-intellectual}

Villani is one of the most recognizable mathematicians in the world,
known for: - His distinctive style (spider brooches, cravats, elaborate
attire) - His appearances in film (he appeared in the documentary
\emph{The Great Mystery of Mathematics}) - His books for the general
public (\emph{Theorems Live!}, \emph{The Mathematical House}) - His TED
talks and public lectures

\begin{center}\rule{0.5\linewidth}{0.5pt}\end{center}

\section{18. Awards and Recognition}\label{ch2010cedricvillani:awards-and-recognition}

\textbf{2010:} - Fields Medal (International Mathematical Union)

\textbf{2009:} - Fermat Prize (Institut de Mathématiques de Toulouse) -
Henri Poincaré Prize (International Association of Mathematical Physics)

\textbf{2008:} - Prize of the Jacques Herbrand Foundation (Académie des
Sciences)

\textbf{2007:} - European Mathematical Society Prize

\textbf{2006:} - Prix Charles Lagrange (Académie Royale de Belgique)

\textbf{2004:} - Louis Armand Prize (Académie des Sciences)

\textbf{2003:} - Prix Jacques Neveu (Académie des Sciences)

\textbf{Memberships:} - Corresponding Member, French Academy of Sciences
- Member, European Academy of Sciences

\begin{center}\rule{0.5\linewidth}{0.5pt}\end{center}

\section{19. Selected Major
Publications}\label{ch2010cedricvillani:selected-major-publications}

\begin{enumerate}
\def\labelenumi{\arabic{enumi}.}
\item
 \textbf{``Landau damping''} (with C. Mouhot). \emph{Acta Mathematica}
 207(1), 29--201, 2011. --- \emph{The proof of nonlinear Landau
 damping.}
\item
 \textbf{``On the rate of convergence to equilibrium for the spatially
 homogeneous Boltzmann equation''} (with L. Desvillettes). \emph{Annals
 of Mathematics} 159(1), 281--340, 2004. --- \emph{Exponential
 convergence for the Boltzmann equation.}
\item
 \textbf{``On the trend to equilibrium for the Fokker--Planck equation
 and the Landau equation.''} \emph{Journal of Statistical Physics}
 124(1), 37--73, 2006.
\item
 \textbf{``Hypocoercivity.''} \emph{Memoirs of the American
 Mathematical Society} 202(950), 2009. --- \emph{Systematic framework
 for hypocoercivity.}
\item
 \textbf{``On the trend to equilibrium for the spatially homogeneous
 Boltzmann equation.''} \emph{Inventiones Mathematicae} 149(2),
 285--316, 2002.
\item
 \textbf{``A review of mathematical topics in collisional kinetic
 theory''} (with L. Desvillettes). \emph{Handbook of Mathematical Fluid
 Dynamics} 1, 2002.
\item
 \textbf{``Optimal transport, old and new.''} \emph{Grundlehren der
 mathematischen Wissenschaften} 338, Springer, 2009. --- \emph{The
 definitive reference on optimal transport.}
\item
 \textbf{``Topics in optimal transportation.''} \emph{Graduate Studies
 in Mathematics} 58, AMS, 2003. --- \emph{Graduate-level introduction
 to optimal transport.}
\item
 \textbf{``Regularity of optimal transport maps''} (with G. De
 Philippis and A. Figalli). \emph{Journal of the European Mathematical
 Society} 16(6), 2014.
\item
 \textbf{``From the Boltzmann equation to the incompressible
 Navier--Stokes equations''} (with L. Saint-Raymond).
 \emph{Communications on Pure and Applied Mathematics} 57(1), 2004.
\end{enumerate}

\begin{center}\rule{0.5\linewidth}{0.5pt}\end{center}

\section{20. Ongoing Work and Future
Directions}\label{ch2010cedricvillani:ongoing-work-and-future-directions}

\subsection{20.1 Hydrodynamic Limits}\label{ch2010cedricvillani:hydrodynamic-limits}

The derivation of fluid equations (Euler, Navier--Stokes) from kinetic
equations remains an active area. Villani continues to work on: - The
Boltzmann-to-Navier-Stokes limit - The Vlasov-to-magnetohydrodynamics
limit - Quantitative error estimates for these limits

\subsection{20.2 Biological Applications}\label{ch2010cedricvillani:biological-applications}

Villani has recently been applying kinetic theory to biological
problems: - \textbf{Chemotaxis:} The movement of organisms in response
to chemical gradients - \textbf{Population dynamics:} Kinetic models of
interacting populations - \textbf{Epidemiology:} Kinetic models of
disease spread

\subsection{20.3 Machine Learning}\label{ch2010cedricvillani:machine-learning}

Villani's report on AI for the French government led to an interest in
the mathematical foundations of machine learning. He has been working
on: - \textbf{Optimal transport in machine learning:} Using Wasserstein
distances for generative models - \textbf{Convergence of optimization
algorithms:} Gradient flows in Wasserstein space - \textbf{The
mathematics of neural networks:} Understanding the approximation
capabilities of deep networks

\subsection{20.4 Public Policy}\label{ch2010cedricvillani:public-policy}

While no longer in parliament, Villani remains active in science policy,
advising on: - AI regulation - Research funding - Mathematics education
- International scientific collaboration

\begin{center}\rule{0.5\linewidth}{0.5pt}\end{center}

\section{21. Conclusion}\label{ch2010cedricvillani:conclusion}

Cédric Villani received the Fields Medal in 2010 for two monumental
achievements: the proof of exponential convergence to equilibrium for
the Boltzmann equation, and the proof of nonlinear Landau damping for
the Vlasov--Poisson equation. Both results resolved problems that had
been open for over half a century and transformed our understanding of
kinetic theory.

Beyond these specific results, Villani's contributions to optimal
transport --- both through his research and through his two monumental
books --- have had a substantial impact on mathematics. Optimal
transport, once a specialized topic, is now a central tool in analysis,
geometry, probability, and applied mathematics.

Villani's career also demonstrates the importance of mathematical
exposition and public engagement. His books, lectures, and public
service have made mathematics accessible to a wide audience and have
inspired a new generation of mathematicians.

The Fields Medal citation --- \emph{``for his proofs of nonlinear Landau
damping and convergence to equilibrium for the Boltzmann equation''} ---
captures his two greatest achievements. But his legacy extends far
beyond these results, encompassing the development of optimal transport
as a central mathematical tool, the creation of the hypocoercivity
framework, and his role as a public intellectual bringing mathematics to
society.

\begin{center}\rule{0.5\linewidth}{0.5pt}\end{center}

\begin{center}\rule{0.5\linewidth}{0.5pt}\end{center}

\section*{Exercises}
\addcontentsline{toc}{section}{Exercises}

\begin{exercise} Show that the Wasserstein distance $W_1(\\mu, \\nu)$ metrizes weak convergence of probability measures on a compact space. \end{exercise}

\begin{exercise} Verify the Brenier theorem: the optimal transport map from a Gaussian to any other measure is the gradient of a convex function. \end{exercise}

\begin{exercise} Compute the entropy production for the Boltzmann equation in the relaxation-time approximation. \end{exercise}

\begin{exercise} Prove the Talagrand transportation inequality for the Gaussian measure. \end{exercise}

\begin{center}\rule{0.5\linewidth}{0.5pt}\end{center}

\section*{Summary}
\addcontentsline{toc}{section}{Summary}

Villani developed the theory of optimal transport and its applications to PDEs, Riemannian geometry, and mathematical physics. His work on Landau damping in plasma physics and the entropy method for dissipation inequalities established deep connections between analysis, probability, and geometry.

\section{22. References and Further
Reading}\label{ch2010cedricvillani:references-and-further-reading}

\subsection{Official Sources}\label{ch2010cedricvillani:official-sources}

\begin{itemize}
\tightlist
\item
 ICM 2010 Fields Medal Announcement:
 https://www.icm2010.in/prize-winners-2010/fields-medal-cedric-villani
\item
 IMU Fields Medal Citations:
 https://www.mathunion.org/imu-awards/fields-medal/fields-medal-2010
\end{itemize}

\subsection{Biographical and Overview
Articles}\label{ch2010cedricvillani:biographical-and-overview-articles}

\begin{itemize}
\tightlist
\item
 Villani, C. ``The Crowded Journey of a Mathematician.'' \emph{Notices
 of the AMS} 58(1), 2011.
\item
 Ambrosio, L. ``The Work of Cédric Villani.'' ICM 2010 Laudatio.
\item
 Villani, C. \emph{Birth of a Theorem: A Mathematical Adventure.}
 Farrar, Straus and Giroux, 2015.
\end{itemize}

\subsection{Key Technical References}\label{ch2010cedricvillani:key-technical-references}

\begin{itemize}
\tightlist
\item
 Mouhot, C. and Villani, C. ``Landau damping.'' \emph{Acta Mathematica}
 207 (2011), 29--201.
\item
 Desvillettes, L. and Villani, C. ``On the rate of convergence to
 equilibrium for the spatially homogeneous Boltzmann equation.''
 \emph{Annals of Mathematics} 159 (2004), 281--340.
\item
 Villani, C. ``Hypocoercivity.'' \emph{Memoirs of the AMS} 202 (2009).
\item
 Villani, C. \emph{Optimal Transport: Old and New.} Springer, 2009.
\end{itemize}

\subsection{Background Reading}\label{ch2010cedricvillani:background-reading}

\begin{itemize}
\tightlist
\item
 Boltzmann, L. ``Weitere Studien über das Wärmegleichgewicht unter
 Gasmolekülen.'' \emph{Wiener Berichte} 66 (1872), 275--370.
\item
 Landau, L. ``On the vibration of the electronic plasma.''
 \emph{Journal of Physics USSR} 10 (1946), 25--34.
\item
 Cercignani, C. \emph{The Boltzmann Equation and Its Applications.}
 Springer, 1988.
\item
 Ambrosio, L., Gigli, N., and Savaré, G. \emph{Gradient Flows in Metric
 Spaces and in the Space of Probability Measures.} Birkhäuser, 2005.
\end{itemize}

\begin{center}\rule{0.5\linewidth}{0.5pt}\end{center}

\section*{See Also}
\addcontentsline{toc}{section}{See Also}

This chapter relates to the following work by other Fields Medalists:
\begin{itemize}
  \item \textbf{\hyperref[ch1994pierrelouislions:pierre-louis-lions-fields-medal-1994-a-comprehensive-report-on-his-work-and-impact]{See Chapter on related work}} --- Villani is Lions' student.
  \item \textbf{\hyperref[ch2006terencetao:terence-tao-fields-medal-2006-a-comprehensive-report-on-his-work-and-impact]{See Chapter on related work}} --- Tao's work on PDEs relates to Villani.
\end{itemize}

